% ============================================================
% Submission LaTeX source — International Business Review (Elsevier)
% Document class for final submission: Elsevier elsarticle.cls (review,12pt) — Elsevier; article fallback để xác minh
% Compile: xelatex 04_manuscript_latex.tex (x2). Figures submitted separately.
% ============================================================
% Options for packages loaded elsewhere
\PassOptionsToPackage{unicode}{hyperref}
\PassOptionsToPackage{hyphens}{url}
\PassOptionsToPackage{dvipsnames,svgnames,x11names}{xcolor}
%
\documentclass[
  12pt,a4paper]{article}
\usepackage{amsmath,amssymb}
\usepackage{setspace}
\usepackage{iftex}
\ifPDFTeX
  \usepackage[T1]{fontenc}
  \usepackage[utf8]{inputenc}
  \usepackage{textcomp} % provide euro and other symbols
\else % if luatex or xetex
  \usepackage{unicode-math} % this also loads fontspec
  \defaultfontfeatures{Scale=MatchLowercase}
  \defaultfontfeatures[\rmfamily]{Ligatures=TeX,Scale=1}
\fi
\usepackage{lmodern}
\ifPDFTeX\else
  % xetex/luatex font selection
  \setmainfont[]{TeX Gyre Termes}
  \setmathfont[]{TeX Gyre Termes Math}
\fi
% Use upquote if available, for straight quotes in verbatim environments
\IfFileExists{upquote.sty}{\usepackage{upquote}}{}
\IfFileExists{microtype.sty}{% use microtype if available
  \usepackage[]{microtype}
  \UseMicrotypeSet[protrusion]{basicmath} % disable protrusion for tt fonts
}{}
\makeatletter
\@ifundefined{KOMAClassName}{% if non-KOMA class
  \IfFileExists{parskip.sty}{%
    \usepackage{parskip}
  }{% else
    \setlength{\parindent}{0pt}
    \setlength{\parskip}{6pt plus 2pt minus 1pt}}
}{% if KOMA class
  \KOMAoptions{parskip=half}}
\makeatother
\usepackage{xcolor}
\usepackage[margin=2.5cm]{geometry}
\usepackage{longtable,booktabs,array}
\usepackage{calc} % for calculating minipage widths
% Correct order of tables after \paragraph or \subparagraph
\usepackage{etoolbox}
\makeatletter
\patchcmd\longtable{\par}{\if@noskipsec\mbox{}\fi\par}{}{}
\makeatother
% Allow footnotes in longtable head/foot
\IfFileExists{footnotehyper.sty}{\usepackage{footnotehyper}}{\usepackage{footnote}}
\makesavenoteenv{longtable}
\setlength{\emergencystretch}{3em} % prevent overfull lines
\providecommand{\tightlist}{%
  \setlength{\itemsep}{0pt}\setlength{\parskip}{0pt}}
\setcounter{secnumdepth}{-\maxdimen} % remove section numbering
\ifLuaTeX
\usepackage[bidi=basic]{babel}
\else
\usepackage[bidi=default]{babel}
\fi
\babelprovide[main,import]{english}
\ifPDFTeX
\else
\babelfont{rm}[]{TeX Gyre Termes}
\fi
% get rid of language-specific shorthands (see #6817):
\let\LanguageShortHands\languageshorthands
\def\languageshorthands#1{}

\setcounter{secnumdepth}{-1}  % giữ numbering thủ công trong manuscript
\usepackage{booktabs}
\usepackage{longtable}
\usepackage{newunicodechar}
\newunicodechar{₀}{\textsubscript{0}}\newunicodechar{₁}{\textsubscript{1}}
\newunicodechar{₂}{\textsubscript{2}}\newunicodechar{₃}{\textsubscript{3}}
\newunicodechar{₄}{\textsubscript{4}}\newunicodechar{₅}{\textsubscript{5}}
\newunicodechar{≈}{$\approx$}\newunicodechar{≤}{$\leq$}\newunicodechar{≥}{$\geq$}
\newunicodechar{×}{$\times$}\newunicodechar{−}{$-$}
\ifLuaTeX
  \usepackage{selnolig}  % disable illegal ligatures
\fi
\IfFileExists{bookmark.sty}{\usepackage{bookmark}}{\usepackage{hyperref}}
\IfFileExists{xurl.sty}{\usepackage{xurl}}{} % add URL line breaks if available
\urlstyle{same}
\hypersetup{
  pdftitle={Internationalization and Firm Performance across 50 Asian and Pacific Economies: Institutional Regimes, Digital Capabilities, and Managerial Characteristics as Contingency Factors},
  pdflang={en},
  colorlinks=true,
  linkcolor={black},
  filecolor={Maroon},
  citecolor={black},
  urlcolor={blue},
  pdfcreator={LaTeX via pandoc}}

\title{Internationalization and Firm Performance across 50 Asian and
Pacific Economies: Institutional Regimes, Digital Capabilities, and
Managerial Characteristics as Contingency Factors}
\author{}
\date{}

\begin{document}
\maketitle
\begin{abstract}
We examine whether the internationalization--performance (I--P)
relationship holds across a large and institutionally diverse set of
Asian and Pacific economies and identify the conditions under which
firms convert export intensity into productive performance. Using
microdata from 88,869 firms across 50 economies and 103 country-year
waves of the World Bank Enterprise Survey (2003--2025; including Japan's
inaugural 2025 wave), spanning all six Institutional Context Regime
Variation (ICRV) groups, we estimate a hierarchical sequence of models
on currency-neutral labour productivity (standardised within
economy-year) with two-way (economy and year) fixed effects throughout
and standard errors clustered by economy. We find robust support for an
inverted-U I--P relationship: the quadratic baseline (M2, N=81,022)
yields a turning point of 51.5\% of export sales, tightening to 43.6\%
once technological capability and digital adoption are controlled (M4,
N=79,080; Lind--Mehlum p \textless{} .001 throughout). Technological
capability (TCI, +0.108) and digital adoption (DAI, +0.219) both raise
the performance level, but neither significantly reshapes the pooled
curve; both are level-shifters, not curve-reshapers. The curve structure
is instead governed by institutional regime: estimating the turning
point within each ICRV group reveals a three-zone pattern: a sharp
inverted-U in the lower-middle transition regime (TP ≈ 43\%), a
near-linear relationship in the strongest-institution regime (turning
point beyond the observable range, ≈ 80\%), and dissolution of the
inverted-U in the weakest regimes (Emerging ≈ 37\%; the Pacific SIDS
group shows no inverted-U). Female top management carries a negative
coefficient (−0.104, p \textless{} .001). These findings offer a
parsimonious resolution to Asian I--P heterogeneity: the relationship is
inverted-U in form, but it is institutional regime, not capability per
se, that determines whether the inverted-U is sharp, straightened, or
absent.
\end{abstract}

\setstretch{2.0}
\textbf{Keywords:} internationalization--performance; inverted-U;
institutional regime; digital capability; Asia-Pacific; World Bank
Enterprise Surveys

\hypertarget{introduction}{%
\subsection{1. Introduction}\label{introduction}}

A central question in international business research concerns whether
and how export internationalization translates into firm performance (Lu
\& Beamish, 2004; Johanson \& Vahlne, 2009). Meta-analyses document a
mean correlation of approximately r = .07 across several thousand
studies (Bausch \& Krist, 2007; Kirca et al., 2012; Marano et al.,
2016), yet the variance across studies vastly exceeds what sampling
error can explain. In Asia (a region that accounts for over half of
global manufacturing exports; UNCTAD, 2023), the evidence is especially
fragmented: prior studies report positive relationships (Pangarkar,
2008; Korea), inverted-U curves (Chiao et al., 2006; Taiwan), or null
and negative effects (Mondal et al., 2022; India). The standard
explanation invokes context as a moderator (Marano et al., 2016; Kirca
et al., 2012), but empirical tests simultaneously across institutionally
diverse Asian economies remain rare.

This paper addresses that gap. We pool 88,869 firm-observations from 50
Asian and Pacific economies spanning 2003--2025 (including Japan's
inaugural 2025 wave), harmonised from the World Bank Enterprise Survey
(WBES). Our sample spans all six ICRV regime groups, from advanced
innovation-driven economies (Japan, Singapore, Korea, Taiwan, HongKong,
Israel) to Pacific Small Island Developing States (Vanuatu, Solomon
Islands, Fiji, Tonga, Samoa, Kiribati, Papua New Guinea) and Gulf
advanced-resource economies (Bahrain, Kuwait, Qatar, Brunei, Saudi
Arabia, Oman), all classified through the Institutional Context Regime
Variation (ICRV) framework that anchors the broader dissertation
programme of which this study is part.

Our contributions are threefold. First, we provide the largest
WBES-based test of the inverted-U I--P hypothesis in Asia, confirming a
capability-adjusted turning point of 43.6\% foreign sales intensity
(51.5\% in the unconditioned quadratic; Lind--Mehlum p \textless{} .001
throughout) on a harmonised 50-economy frame. Second, we show that both
technological capability (TCI, +0.108) and digital adoption (DAI,
+0.219) are strong performance \textbf{level-shifters} but that neither
significantly reshapes the pooled I--P curve, a finding that disciplines
capability-centred accounts of internationalization performance. Third,
and most consequentially, we show that the \textbf{shape} of the curve
is governed by institutional regime, and not as a smooth gradient but as
a three-zone structure: a sharp inverted-U in transition economies
(turning point ≈ 43\%), near-linearity in the strongest-institution
economies (turning point beyond the observable range), and dissolution
(reversing into a Forced Internationalization Penalty among the Pacific
SIDS) in the weakest regimes.

The remainder of the paper is organised as follows. Section 2 develops
the theoretical framework and formal hypotheses. Section 3 describes
data and methods. Section 4 presents results. Section 5 discusses
implications and limitations. Section 6 concludes.

\begin{center}\rule{0.5\linewidth}{0.5pt}\end{center}

\hypertarget{theoretical-framework-and-hypotheses}{%
\subsection{2. Theoretical Framework and
Hypotheses}\label{theoretical-framework-and-hypotheses}}

\hypertarget{theoretical-foundations}{%
\subsubsection{2.1 Theoretical
Foundations}\label{theoretical-foundations}}

We draw on four theoretical traditions. The Uppsala model (Johanson \&
Vahlne, 1977, 2009) describes internationalization as a sequential
process in which firms accumulate experiential knowledge, progressively
reducing psychic distance and coordination costs. The Resource-Based
View (Barney, 1991; Wernerfelt, 1984) positions firm-level capabilities
(technological, digital, and managerial) as the mediating mechanism
linking export exposure to performance outcomes. Institutional Theory
(North, 1990; Scott, 2008) directs attention to how country-level formal
and informal institutions shape the cost structure of cross-border
transactions, and how home-country uncertainty conditions the
internationalization--performance link (Cuervo-Cazurra, Ciravegna,
Melgarejo, \& Lopez, 2018). Upper Echelons Theory (Hambrick \& Mason,
1984; Hambrick, 2007) emphasises that top management team attributes
moderate how firms respond to strategic opportunities and constraints.

\hypertarget{the-non-linear-ip-relationship-h1}{%
\subsubsection{2.2 The Non-Linear I--P Relationship
(H1)}\label{the-non-linear-ip-relationship-h1}}

Contractor et al.~(2003) formalised the three-stage S-curve hypothesis:
early internationalization carries fixed entry costs that depress
performance; intermediate internationalization yields learning and scale
economies; high internationalization may again impose coordination and
agency costs. Empirically, however, the data more consistently support
the inverted-U: a special case in which the third stage does not emerge
at the export intensities typically observed in developing-economy WBES
data (Gomes \& Ramaswamy, 1999; Lu \& Beamish, 2004).

In the Asian context, three country-specific studies in our broader
dissertation programme provide prior evidence. For Vietnam (P3), the
inverted-U is confirmed with a turning point at 39--46\% using a
WBES-based instrumental variable design. For China (P5), an inverted-U
turning point of 47--49\% is replicated across specifications with
Paternoster cross-cohort tests (p = .545). For Singapore (P4), the
relationship is predominantly positive with a very late turning point
(\textasciitilde88.6\%), consistent with near-monotonic returns in a
high-institutional, digitally saturated economy. Pooled across a much
larger and institutionally heterogeneous sample, we expect the
inverted-U to emerge at an intermediate turning point.

\begin{quote}
\textbf{H1:} The relationship between export intensity (FSTS) and firm
labour productivity (ln LP) is an inverted-U, with a statistically
confirmed turning point at intermediate export intensities, in a pooled
sample of Asian and Pacific firms.
\end{quote}

\hypertarget{technological-capability-as-a-curve-amplifier-h2}{%
\subsubsection{2.3 Technological Capability as a Curve Amplifier
(H2)}\label{technological-capability-as-a-curve-amplifier-h2}}

Technological capability, operationalised as a composite of quality
certification, foreign-technology adoption, process innovation, and
R\&D, raises absorptive capacity (Cohen \& Levinthal, 1990), enabling
firms to extract greater performance value per unit of export exposure.
Critically, we expect TCI to operate as a \emph{curve amplifier}: at low
export intensities, high-TCI firms gain more per unit of
internationalization (steeper ascending limb); at high intensities, they
experience sharper diminishing returns as coordination demands exceed
even their superior capabilities (steeper descending limb). The net
effect is a shift and steepening of the inverted-U, not just a parallel
upward shift.

This prediction is consistent with P3 Vietnam (β(TCI) = +0.179,
IV-identified) and P5 China (β(TCI) increasing from +0.28 to +0.43
across cohorts), both showing TCI as a positive level shifter. P4
Singapore shows negligible TCI moderation of the curve, likely because
near-universal adoption in a high-ICRV economy leaves insufficient
variance for moderation.

\begin{quote}
\textbf{H2:} TCI amplifies the I--P curve, raising the average
performance level and steepening both the ascending and descending limbs
of the inverted-U.
\end{quote}

\hypertarget{digital-adoption-as-a-level-enhancer-h3}{%
\subsubsection{2.4 Digital Adoption as a Level Enhancer
(H3)}\label{digital-adoption-as-a-level-enhancer-h3}}

Digital adoption, measured through website presence and
electronic-payment share, has been associated with reduced communication
and transaction costs (Verhoef et al., 2021). In the CDCM
(Context-Contingent Digital Capability Model) developed in the
dissertation, DAI's effect depends on the interaction of regime quality,
FSTS level, and digital market saturation (Authors, 2024). Specifically,
in institutionally advanced economies where digital infrastructure is
mature, DAI provides universal performance lifts (as in P4 Singapore:
β(DAI) positive and significant). In emerging-economy contexts with
substantial selection into digital adoption (P3 Vietnam: IV-estimated
β(DAI) ≈ 0, consistent with selection bias in OLS), DAI's curve-shaping
role is attenuated. Across the pooled multi-economy sample, we expect
DAI to raise performance levels broadly across the sample (level
effect). In a fully specified model that also accounts for managerial
quality (which correlates with digital adoption), the residual DAI
variation may reveal a capability-complementarity curve effect:
digitally equipped firms at low FSTS face sunk adoption costs that
compress near-term productivity, while digital infrastructure supports
coordination efficiency at high FSTS where cross-border management
demands are greatest.

\begin{quote}
\textbf{H3:} DAI has a positive main effect on labour productivity and,
when managerial characteristics are controlled, reshapes the I--P curve
by compressing the ascending limb at low export intensities and
attenuating performance decline at high export intensities.
\end{quote}

\hypertarget{managerial-characteristics-h4}{%
\subsubsection{2.5 Managerial Characteristics
(H4)}\label{managerial-characteristics-h4}}

Upper Echelons Theory predicts that top manager experience and gender
composition affect strategic cognition and decision-making (Hambrick \&
Mason, 1984). Manager experience reduces coordination errors in
cross-border operations; female top management is associated with
stakeholder orientation, relationship capital, and risk management, all
relevant in institutionally heterogeneous export markets (Richard et
al., 2019). We expect both to raise the performance level for any given
FSTS. However, given the breadth of the multi-country sample, firm-level
managerial differences are unlikely to systematically reshape the
aggregate FSTS--performance curve, whose shape is primarily determined
by economy-wide institutional and market-maturity factors.

\begin{quote}
\textbf{H4:} Top manager experience and female top manager status have
positive main effects on firm performance, but do not significantly
moderate the FSTS--performance curve shape.
\end{quote}

\hypertarget{institutional-regime-as-a-turning-point-shifter-h5}{%
\subsubsection{2.6 Institutional Regime as a Turning-Point Shifter
(H5)}\label{institutional-regime-as-a-turning-point-shifter-h5}}

North (1990) established that formal institutions reduce transaction
costs; stronger institutions lower the cost of cross-border coordination
and thus shift the performance-maximising export intensity (the
inverted-U turning point). In higher-quality institutional environments,
firms need a lower export share to recoup internationalization costs, so
the turning point occurs earlier. Equivalently, holding FSTS constant,
firms in stronger institutional regimes earn higher returns per unit of
internationalization.

We operationalise institutional regime through the ICRV classification
system developed in the dissertation: six groups ranging from Group I
(Advanced innovation-driven: Singapore, Korea, Taiwan, Israel, Cyprus)
to Group VI (SIDS and small open economies: Vanuatu, Solomon Islands,
Timor-Leste). Higher ICRV group numbers denote weaker institutions. We
expect positive FSTS×ICRV and negative FSTS²×ICRV interactions in Group
III--VI compared to Group I (stronger institutions = higher Group I
coefficient), which is observationally equivalent to the turning point
occurring later in lower-quality institutional environments.

\begin{quote}
\textbf{H5:} The ICRV institutional regime positively moderates the I--P
relationship: firms in higher-quality institutional environments (lower
ICRV group number) achieve the performance-maximising export intensity
at lower FSTS levels.
\end{quote}

\hypertarget{regime-contingent-digital-effects-h6}{%
\subsubsection{2.7 Regime-Contingent Digital Effects
(H6)}\label{regime-contingent-digital-effects-h6}}

The Context-Contingent Digital Capability Model (CDCM) developed in this
dissertation programme predicts that the I--P curve-reshaping role of
digital adoption is not uniform across institutional environments. In
institutionally advanced economies (ICRV Group I--II), digital
infrastructure is near-universal and close to the competitive baseline;
the marginal performance differentiation from website and e-payment
adoption is accordingly lower. In institutionally weaker economies (ICRV
Group IV--VI), digital adoption remains differentiating: digital market
access tools lower the liability-of-foreignness cost and provide
cross-border coordination support that institutional voids cannot
substitute. This motivates examining whether any curve-shaping role of
digital capability (predicted by H3) is concentrated in
weaker-institution economies rather than uniform across the pool: a
prediction we assess through the within-regime turning-point analysis
(Table 4), since the reconciled pooled models find no significant DAI
curve-reshaping (H3 level effect confirmed, curvature NS).

This prediction is consistent with the digital-complementarity strand of
institutional theory (North, 1990; Stallkamp \& Schotter, 2021): where
formal institutions provide reliable contract enforcement, IP
protection, and market information, digital infrastructure adds
incrementally; where institutional voids prevail, digital tools provide
a non-trivial substitute for missing formal mechanisms. H6 therefore
does not predict stronger \emph{level} effects in weak-institution
contexts (which could reflect many confounds), but specifically predicts
a stronger \emph{curve-shaping} interaction between digital adoption and
export intensity: a conditional complementarity claim.

\begin{quote}
\textbf{H6:} The ICRV institutional regime moderates the
digital-adoption curve-reshaping effect: digital adoption's compressive
effect on the I--P ascending limb and its buffering effect at high FSTS
(H3) are stronger in weaker institutional environments (higher ICRV
group number), where digital infrastructure provides greater competitive
differentiation in cross-border market access.
\end{quote}

\begin{center}\rule{0.5\linewidth}{0.5pt}\end{center}

\hypertarget{data-and-methods}{%
\subsection{3. Data and Methods}\label{data-and-methods}}

\hypertarget{data-source}{%
\subsubsection{3.1 Data Source}\label{data-source}}

Data derive from the World Bank Enterprise Surveys (WBES; World Bank,
2025), a stratified random-sample survey of formal private-sector
enterprises conducted periodically across over 130 countries. We use
microdata from 50 Asian and Pacific economies across 103 country-year
waves spanning 2003--2025, including Japan's inaugural WBES wave (2025).
This coverage encompasses six UN M.49 Asian sub-regions: East Asia
(China, HongKong, Korea, Mongolia, Taiwan), Southeast Asia (Brunei,
Cambodia, Indonesia, Laos, Malaysia, Myanmar, Philippines, Singapore,
Thailand, Timor-Leste, Vietnam), South Asia (Afghanistan, Bangladesh,
Bhutan, India, Maldives, Nepal, Pakistan, Sri Lanka), Central Asia
(Kazakhstan, Kyrgyz Republic, Tajikistan, Turkmenistan, Uzbekistan),
West Asia (Armenia, Bahrain, Cyprus, Iraq, Israel, Jordan, Kuwait,
Lebanon, Oman, Qatar, Saudi Arabia, Yemen), and Pacific/SIDS (Fiji,
Kiribati, Papua New Guinea, Samoa, Solomon Islands, Tonga, Vanuatu).

The final analytical frame contains 88,869 firm-observations across 50
economies and 103 country-year pairs; the main regression sample
(outcome and FSTS non-missing) is 81,022. Some early waves predate the
digital-capability and manager variables used in the moderation models;
they contribute to the baseline models but drop from later
specifications due to item missingness. WBES designs are cross-sectional
within each wave, with standardised instruments ensuring common variable
definitions across economies and years. We apply three harmonisation
steps: (1) variable-priority mapping for candidate items across PICS3
(2009--2013), Standardised (2014--2018), and BREADY/BEE (2019--2025)
questionnaire generations; (2) encoding-robust reading (primary UTF-8
with Latin-1/CP1252 fallback); (3) deduplication, selecting the largest
file per (country, year) pair when multiple uploads exist.

\hypertarget{variables}{%
\subsubsection{3.2 Variables}\label{variables}}

\textbf{Dependent variable.} Labour productivity is log annual sales per
permanent worker, ln(LP) = ln(d2 / l1), winsorised at the 1st/99th
percentiles within each survey wave and then \textbf{standardised within
each economy-year cell} (lp\_z, mean 0 / SD 1). Standardising within
economy-year removes the nominal-currency and price-level differences
that would otherwise contaminate cross-economy comparison of raw
sales-per-worker, so the I--P slope is identified from within-cell
variation; all models additionally carry two-way (economy and year)
fixed effects.

\textbf{Foreign sales intensity (FSTS).} Defined as the share of sales
from exports, both direct and indirect, (d3b + d3c) / 100, ranging from
0 to 1. We mean-centre FSTS at the sample mean for quadratic and
interaction models (fsts\_c); the quadratic term fsts\_c² is computed
from the centred value. The sample mean is 9.0\%.

\textbf{Technological capability index (TCI\_z).} A two-item composite
of internationally recognised quality certification (b8: has ISO or
other certification) and foreign-technology adoption (e6: uses
foreign-licensed technology), standardised to mean 0, SD 1 within
economy-year.

\textbf{Digital adoption index (DAI).} Website presence (c22b), a binary
Tier-1 indicator of digital participation.

\textbf{Manager characteristics.} Top manager experience in sector (b7,
in years) and a binary indicator for female top manager (b7a / b6a).

\textbf{Controls.} Female ownership share (b4 \textgreater{} 0.5:
majority female-owned, binary), foreign ownership percentage (b6a /
100), firm age (in years, current year − b5), and log permanent
employees (ln l1\_workers) as a size proxy.

\textbf{Institutional regime (ICRV group).} Integer 1--6 based on the
ICRV classification; used as a continuous moderator in M8 (its main
effect absorbed by economy fixed effects) and as the grouping variable
for the within-regime turning-point analysis (Table 4).

\hypertarget{estimation-strategy}{%
\subsubsection{3.3 Estimation Strategy}\label{estimation-strategy}}

We estimate a hierarchical sequence of eight models (Table 2). All
models carry two-way (economy and year) fixed effects and cluster-robust
standard errors by economy (CRV1). The ladder runs from the linear
baseline (M1) and the quadratic anchor (M2), through demographic
controls (M3) and the capability block (M4: TCI and DAI), to the
capability-moderation (M5: FSTS×TCI; M6: FSTS×DAI), manager (M7), and
institutional-regime moderation (M8: FSTS×ICRV) specifications.

The inverted-U shape is formally confirmed using the Lind--Mehlum (2010)
test, which requires: (a) β₁ \textgreater{} 0, β₂ \textless{} 0, and (b)
the turning point lies strictly within the observed FSTS range, and (c)
the slope at the extremes satisfies the directional prediction. We
report the turning point as: TP = −β₁ / (2β₂) + mean(FSTS), where β₁ and
β₂ are estimated on the mean-centred FSTS\_c scale and the result is
expressed as a proportion (0--1) of total sales.

For moderation tests (M5--M8), we include interaction terms
FSTS×moderator and FSTS²×moderator and report the individual coefficient
p-values, following standard practice for specifying and interpreting
interaction effects in regression (Aiken \& West, 1991; Dawson, 2014).

\hypertarget{sample-sizes}{%
\subsubsection{3.4 Sample Sizes}\label{sample-sizes}}

The analytic frame is 88,869 firms across 50 economies. The quadratic
anchor (M2: outcome + FSTS non-missing) is N = 81,022. Adding the
capability block (M4: TCI and DAI) gives N = 79,080; demographic
controls (M3) N = 78,970; the manager models (M7) drop to N = 75,029
because manager experience and gender are missing in some waves. The
modest attrition across the ladder reflects item missingness in the
moderator and manager variables, which is more prevalent in smaller
economies and early WBES waves. All N's are reproducible via
\texttt{scripts/p7\_full\_ladder.py}.

\begin{center}\rule{0.5\linewidth}{0.5pt}\end{center}

\hypertarget{results}{%
\subsection{4. Results}\label{results}}

\hypertarget{descriptive-statistics-and-correlations}{%
\subsubsection{4.1 Descriptive Statistics and
Correlations}\label{descriptive-statistics-and-correlations}}

The full analytic sample spans 50 economies and 103 country-year waves.
The mean FSTS across firms is approximately 9\% (median ≈ 0), reflecting
the fact that most WBES firms are domestically oriented; exporters (FSTS
\textgreater{} 0) represent roughly one in five firms. Because the
dependent variable is standardised within each economy-year cell,
productivity is directly comparable within the estimation regardless of
national currency or price level. Pairwise correlations among focal
variables are modest, and variance inflation factors in the full model
are below 3.5, ruling out multicollinearity concerns.

\hypertarget{main-effect-inverted-u-h1}{%
\subsubsection{4.2 Main Effect: Inverted-U
(H1)}\label{main-effect-inverted-u-h1}}

\emph{Table 2, Models M1--M4.}

The linear baseline (M1) yields a positive but small FSTS slope (β =
+0.219, p = .002). Adding the quadratic term (M2, N=81,022) yields β₁ =
+1.188 (p \textless{} .001) and β₂ = −1.398 (p \textless{} .001),
consistent with an inverted-U. Following Haans, Pieters, and He (2016),
the formal conditions for a genuine inverted-U are satisfied: (C1) β₁
\textgreater{} 0 (p \textless{} .001), a positive ascending limb; (C2)
β₂ \textless{} 0 (p \textless{} .001), concavity; (C3) the turning point
TP ≈ 51.5\% lies within the observed FSTS range; and (C4) the marginal
effect is positive at the lower boundary and negative at the upper
boundary, with opposite signs confirmed by the Lind--Mehlum (2010) test
(joint p \textless{} .001). The turning point tightens as capability is
controlled: adding demographic controls (M3) gives TP = 43.8\%, and
adding technological capability and digital adoption (M4, N=79,080)
gives TP = 43.6\%, the headline estimate, matching the dissertation's
canonical value. The stability of the turning point in the
low-to-mid-40\% range once capability is held constant indicates that
the inverted-U is not an artefact of capability-based selection into
exporting.

\textbf{H1 is supported.} The I--P relationship is robustly inverted-U
in the pooled sample of 50 Asian and Pacific economies, with a
capability-adjusted turning point of 43.6\% foreign sales intensity
(51.5\% in the unconditioned quadratic).

\hypertarget{technological-capability-moderation-h2}{%
\subsubsection{4.3 Technological Capability Moderation
(H2)}\label{technological-capability-moderation-h2}}

\emph{Table 2/3, Models M4--M5.}

Technological capability enters with a positive and highly significant
level effect (M4: TCI = +0.108, p \textless{} .001): higher-capability
firms are more productive at every level of internationalization. The
interaction model M5 adds FSTS×TCI and FSTS²×TCI; neither is significant
(FSTS×TCI = +0.068, FSTS²×TCI = −0.119, both p \textgreater{} .10), and
the turning point is essentially unchanged (TP = 44.0\%). Capability
raises the floor; it does not bend the curve.

\textbf{H2 is partially supported (level effect confirmed; curvature
effects NS).} TCI is a foundational level-shifter across the 50-economy
pool. The curve-moderation found in single-country studies (P3 Vietnam,
P5 China) does not survive aggregation across the institutionally
diverse pool, consistent with capability's curve-shaping role being
institution-contingent rather than universal.

\hypertarget{digital-adoption-moderation-h3}{%
\subsubsection{4.4 Digital Adoption Moderation
(H3)}\label{digital-adoption-moderation-h3}}

\emph{Table 2/3, Models M4 and M6.}

Digital adoption carries the largest capability level effect in the
model (M4: DAI = +0.219, p \textless{} .001): a substantial productivity
premium for firms with a web presence. The interaction model M6 adds
FSTS×DAI and FSTS²×DAI to test whether digital adoption reshapes the
curve; on this canonical specification, neither interaction is
significant (FSTS×DAI = −0.270, FSTS²×DAI = +0.252, both p
\textgreater{} .10), and the turning point is essentially unchanged (TP
= 47.1\%). The curve-reshaping role attributed to DAI in an earlier
draft does not survive re-estimation on the
within-economy-year-standardised, two-way-fixed-effects specification:
DAI behaves as a level-shifter, not a curve-reshaper, in the pooled
sample.

\textbf{H3 is partially supported (level effect confirmed; curvature
effects NS).} Digital adoption raises firm performance substantially but
does not significantly reshape the pooled I--P curve. The ``digital
shield'' effect (DAI attenuating the descending limb) documented in the
single-country Singapore study (P4) is not detectable once aggregated
across the 50-economy pool, again consistent with the curve-shaping role
of digital capability being institution-contingent rather than
universal.

\hypertarget{managerial-characteristics-h4-1}{%
\subsubsection{4.5 Managerial Characteristics
(H4)}\label{managerial-characteristics-h4-1}}

\emph{Table 2/3, Model M7.}

Adding manager characteristics (M7, N=75,029), manager sectoral
experience is not significant (β = +0.002, p \textgreater{} .10), while
the female-top-manager coefficient is negative and highly significant (β
= −0.104, p \textless{} .001). The inverted-U is unaffected (TP =
44.0\%). The negative female-management coefficient is consistent with
the rest of the dissertation (the single-country Japan study, P10, and
the thesis Chapter 4 likewise estimate a negative female-top-manager
coefficient) and is most plausibly read as compositional selection of
female-led establishments into smaller, lower-margin niches in economies
where female senior leadership remains rare (International Finance
Corporation, 2022; Buchhave et al., 2026), rather than as a leadership
effect. It should not be interpreted causally.

\textbf{H4 is not supported as hypothesised.} Manager experience carries
no significant level effect, and the female-management coefficient is
negative rather than the hypothesised positive. We treat the
female-management coefficient as a descriptive regularity requiring
within-economy, selection-aware analysis, not as evidence on the causal
effect of female leadership.

\hypertarget{institutional-regime-moderation-h5}{%
\subsubsection{4.6 Institutional Regime Moderation
(H5)}\label{institutional-regime-moderation-h5}}

\emph{Table 2/3, Model M8, and the per-ICRV turning-point table.}

Because ICRV regime is time-invariant within an economy, its main effect
is absorbed by the economy fixed effects; only the interaction terms are
identified. In the pooled model M8, neither FSTS×ICRV (β = −0.086) nor
FSTS²×ICRV (β = −0.013) is significant (both p \textgreater{} .10): the
institutional gradient does not manifest as a single linear shift of the
pooled curve. Instead, the institutional structure of the relationship
emerges when the turning point is estimated \textbf{within} each ICRV
group (capability-adjusted M4 form within group):

\begin{itemize}
\tightlist
\item
  \textbf{Lower-middle transition (Group IV, N=42,094): TP ≈ 43\%}, a
  sharp, reliably located inverted-U (β₂ \textless{} 0, p \textless{}
  .001). This is the regime that anchors the pooled estimate.
\item
  \textbf{Strongest-institution regime (Group I Advanced innovation,
  incl.~Japan): TP ≈ 80\%}: the turning point lies near or beyond the
  observable range, so within the data the relationship is effectively
  near-linear and positive.
\item
  \textbf{Weakest regimes (Group V Emerging: TP ≈ 37\%; Group VI Pacific
  SIDS): the inverted-U dissolves}; in the SIDS group the quadratic
  flips sign (no interior maximum), consistent with the Forced
  Internationalization Penalty documented in the companion SIDS study
  (P8).
\end{itemize}

\textbf{H5 is supported, but as existence-moderation rather than
location-moderation.} Institutional regime does not merely shift where a
universal inverted-U peaks; it governs whether the inverted-U is sharp
(transition regimes), straightened (strongest regimes), or absent
(weakest regimes): a three-zone structure.

\hypertarget{synthesis-a-three-zone-institution-contingent-curve}{%
\subsubsection{4.7 Synthesis: A Three-Zone, Institution-Contingent
Curve}\label{synthesis-a-three-zone-institution-contingent-curve}}

Taken together, the moderation results relocate the explanatory weight
from capability to institutions. Capability (TCI, DAI) and managers
shift the \textbf{level} of productivity but do not significantly
reshape the pooled curve. What reshapes the curve is institutional
regime, and it does so not as a smooth gradient but as three discrete
zones. In the middle of the institutional gradient (lower-middle
transition), the inverted-U is sharp and reliably located near a 43\%
turning point; at the top (strongest institutions, including Japan), the
descending limb is pushed beyond the observable range and the
relationship is effectively near-linear; at the bottom (emerging and
Pacific SIDS), the inverted-U dissolves and, in the SIDS group,
reverses. This three-zone reading reconciles the apparently divergent
single-country results in the dissertation (Vietnam and India in the
transition zone show sharp inverted-U; Singapore and Japan at the top
show near-linearity; the SIDS group shows the Forced
Internationalization Penalty) within one harmonised, 50-economy
estimation.

\hypertarget{robustness}{%
\subsubsection{4.8 Robustness}\label{robustness}}

Four supplementary checks bound the inference. \emph{(i) Few-cluster
inference.} Because the per-ICRV-group turning points rest on only 6--16
economies, cluster-robust (CRV1) standard errors can over-reject. A
restricted wild-cluster bootstrap (Cameron, Gelbach, \& Miller, 2008;
Rademacher weights, B = 9,999, clustered by economy) inflates the
per-group curvature p-values: the sharp transition-regime inverted-U
(Group IV, β₂ = −1.012) moves from p = 0.001 (CRV1) to p ≈ 0.10, while
the emerging and SIDS groups remain non-significant. Formal
non-linearity inference is therefore anchored at the pooled level (M2,
49 clusters, β₂ = −1.399, p \textless{} .001) and triangulated with the
companion meta-analysis, with the per-ICRV turning points read as a
descriptive gradient rather than independently powered tests. \emph{(ii)
Multiple testing.} Holm and Bonferroni corrections applied to the
four-term capability-interaction family (FSTS×TCI, FSTS²×TCI, FSTS×DAI,
FSTS²×DAI) leave every term non-significant, so the
capability-as-level-shifter conclusion is not a false negative produced
by uncorrected testing. \emph{(iii) Measurement invariance.} The Tier-1
DAI (website) and TCI component items are populated for
\textasciitilde99\% of firms in all three WBES schema generations; the
DAI mean rises monotonically (0.30 → 0.46 → 0.53), indicating genuine
diffusion rather than a schema discontinuity, so cross-wave comparisons
rest on item-comparable measures. \emph{(iv) Selection into exporting.}
Adding an inverse-Mills term from an export-participation probit to the
controlled specification leaves the curvature and turning point
essentially unchanged (43.6\% → 43.5\%; β₂ = −0.732, p \textless{} .001)
with a non-significant Mills ratio (p = 0.317), indicating no detectable
Melitz-type selection bias in the estimated curvature. Full procedures
and statistics are reported in the online supplement.

\begin{center}\rule{0.5\linewidth}{0.5pt}\end{center}

\hypertarget{discussion}{%
\subsection{5. Discussion}\label{discussion}}

\hypertarget{resolution-of-cross-study-heterogeneity}{%
\subsubsection{5.1 Resolution of Cross-Study
Heterogeneity}\label{resolution-of-cross-study-heterogeneity}}

The central finding (an inverted-U with a capability-adjusted turning
point of 43.6\% foreign sales intensity, robust across 50 economies and
six ICRV groups) offers a parsimonious resolution to apparent
inconsistencies in the Asian I--P literature. Studies finding positive
linear effects (Pangarkar, 2008: Singapore; Cho et al., 2023: Korea) are
working with strong-institution samples whose turning point lies beyond
the observable export-intensity range; their firms are on a long
ascending limb. Studies finding null or negative results may be
capturing weakest-regime samples in which the inverted-U dissolves. Our
within-regime turning-point analysis provides the key refinement: the
relationship is not a single inverted-U with a moving peak but a
\textbf{three-zone} structure: sharp inverted-U in the transition regime
(Group IV, TP ≈ 43\%), near-linearity in the strongest regime (Group I,
TP ≈ 80\%, beyond range), and dissolution in the weakest regimes (Group
V ≈ 37\%; Group VI SIDS, no inverted-U). Whether a single-country study
finds an inverted-U at all therefore depends on which institutional zone
its economy occupies.

\hypertarget{technology-as-amplifier-not-moderator}{%
\subsubsection{5.2 Technology as Amplifier, Not
Moderator}\label{technology-as-amplifier-not-moderator}}

The finding that TCI raises the performance level without reliably
reshaping the I--P curve in the pooled sample distinguishes the
cross-economy result from country-specific evidence. In P3 Vietnam and
P5 China, TCI moderated the curve shape; in the 50-economy pool,
institutional heterogeneity attenuates that moderation signal. This is
consistent with the absorptive capacity argument (Cohen \& Levinthal,
1990): high-TCI firms process international market signals more
effectively, but the curve-shaping benefit requires institutional
complementarity (e.g., strong IP protection, technology transfer
infrastructure) that is inconsistently present across the full ICRV
spectrum. The practical implication is that technological capability
raises the performance floor universally, but its curve-reshaping role
is context-contingent and cannot be assumed at scale.

\hypertarget{digital-adoption-level-effect-not-curve-reshaping}{%
\subsubsection{5.3 Digital Adoption: Level Effect, Not Curve
Reshaping}\label{digital-adoption-level-effect-not-curve-reshaping}}

On the canonical specification, digital adoption is the largest single
capability level-shifter (DAI = +0.219, p \textless{} .001) but does not
significantly reshape the pooled I--P curve. This is a more conservative
and, we argue, more credible reading than an earlier draft that
attributed significant curve-reshaping to DAI: with within-economy-year
standardisation and two-way fixed effects absorbing the price-level and
selection confounds, the apparent curve-reshaping does not survive. The
``digital shield'' mechanism, DAI cushioning the descending limb,
appears in the single-country Singapore study (P4) but is undetectable
in aggregation, indicating that any curve-shaping role of digital
adoption is institution-contingent and operates below the resolution of
a pooled, Tier-1 (website-only) WBES measure. The policy implication is
correspondingly disciplined: digital adoption is a robust lever for
raising the productivity \textbf{level} across the region, but it should
not be sold as a tool for re-shaping the
internationalization-performance trade-off.

\hypertarget{managerial-heterogeneity}{%
\subsubsection{5.4 Managerial
Heterogeneity}\label{managerial-heterogeneity}}

On the canonical specification, manager sectoral experience carries no
significant level effect, and the female-top-manager coefficient is
\textbf{negative} (−0.104, p \textless{} .001) rather than positive.
This sign is consistent across the dissertation: the single-country
Japan study (P10) and the thesis pooled analysis estimate the same
negative coefficient. It cautions against a causal reading. In economies
where female senior leadership remains rare (the regional pool spans
contexts with female-top-manager shares from single digits upward), the
small set of female-led establishments is unlikely to be a random draw;
the most plausible interpretation is compositional selection into
smaller, lower-margin, labour-intensive niches, which a cross-sectional
level regression cannot disentangle from any leadership effect. The null
curve-moderation result (FSTS×manager terms not significant) indicates
that, whatever the source of the level coefficient, manager
characteristics do not help firms navigate the internationalization
trade-off more efficiently. We therefore report the female-management
coefficient as a descriptive regularity requiring within-economy,
selection-aware analysis (e.g., matched samples or panel tracking of
leadership transitions), not as evidence on the causal effect of female
leadership on firm performance.

\hypertarget{limitations}{%
\subsubsection{5.5 Limitations}\label{limitations}}

Three inferential constraints bound these findings. First, WBES is a
cross-sectional design: ``performance'' throughout refers to a
contemporaneous association between FSTS and labour productivity in a
given survey year, not a panel-tracked causal effect. The
instrumental-variable strategy used in the single-country P3 Vietnam
study to address selection into exporting is not available at scale
across 50 economies; the pooled inverse-Mills selection check (Section
4.8) indicates the estimated curvature is not driven by selection into
exporting, but a panel-tracked causal design remains out of reach, and
capability-based selection into exporting is part of what the
capability-adjusted turning point (M4) reveals rather than a nuisance
assumed away. Second, the DAI measure is limited to Tier-1 (website)
capability in the harmonised pool; the richer digital-capability
dimensions captured in more recent BEE-schema waves are not yet
available for the majority of economies, which likely contributes to the
null curve-reshaping result. Third, the dataset covers all 50
ICRV-mapped economies across six regime groups, including Japan's
inaugural 2025 wave; some early waves predate the digital-capability and
manager variables, limiting their contribution to the baseline models.
Future WBES waves would sharpen the moderation estimates, particularly
for the sparser Advanced and SIDS regimes.

\begin{center}\rule{0.5\linewidth}{0.5pt}\end{center}

\hypertarget{conclusions}{%
\subsection{6. Conclusions}\label{conclusions}}

We find that the inverted-U I--P relationship holds robustly across a
50-economy, 103-wave WBES sample (including Japan's inaugural 2025
wave), with a capability-adjusted turning point of 43.6\% foreign sales
intensity. Three conclusions are policy-relevant. First, for firms in
the ascending phase (FSTS \textless{} 43.6\%), intensifying
internationalization is productivity-enhancing; the priority is removing
barriers to export expansion. For firms beyond the turning point, the
constraint is not market access but coordination capacity. Second, both
technological capability and digital adoption raise the performance
level substantially across all FSTS levels without reliably reshaping
the curve; policy should target capability building as a baseline
performance lever rather than as a curve-shifting intervention. Third,
and most important, it is institutional regime, not capability, that
determines the shape of the curve: the inverted-U is sharp in transition
economies (turning point ≈ 43\%), straightens toward near-linearity in
the strongest-institution economies (peak beyond the observable range),
and dissolves in the weakest regimes, reversing into a Forced
Internationalization Penalty in the Pacific SIDS group.
Context-contingent policy must therefore differ by zone, not merely by
where a common peak falls.

The near-term policy context adds urgency to these findings. World Bank
(2026) macro monitoring data for Vietnam, the largest ICRV Group IV
economy in the sample, show manufacturing PMI falling to 45.3
(contraction) in March 2026 as new export orders collapsed following the
opening of a U.S. trade-representative investigation; merchandise
imports had risen 27.8\% year-on-year by the same month, partly driven
by energy supply-chain disruption from the Strait of Hormuz conflict.
These conditions are precisely those under which the right side of the
inverted-U becomes consequential: firms pushed above their optimal FSTS
by earlier export commitments face amplified coordination and input-cost
pressures. Because Vietnam sits in the transition regime where the
inverted-U is sharpest (Group IV turning point ≈ 43\%), Vietnamese
exporters currently above that threshold are among the most exposed to
macro shocks. The capability evidence implies that building
technological and digital capability raises the productivity floor for
these firms, even though, on our pooled estimates, capability does not
by itself flatten the descending limb; the binding margin is the breadth
of capable export participation, not a curve-shifting fix.

Future research should leverage the growing panel dimensions of WBES to
establish causality, and incorporate Tier 3--4 digital capability
measures (AI adoption, cloud computing, platform participation) as these
items become available in the 2025+ WBES waves.

\begin{center}\rule{0.5\linewidth}{0.5pt}\end{center}

\hypertarget{references}{%
\subsection{References}\label{references}}

Aiken, L. S., \& West, S. G. (1991). \emph{Multiple regression: Testing
and interpreting interactions}. Sage.

Barney, J. (1991). Firm resources and sustained competitive advantage.
\emph{Journal of Management, 17}(1), 99--120.
https://doi.org/10.1177/014920639101700108

Buchhave, H., Nguyen, C. V., Vu, C., Nguyen, G. T., \& Zumbyte, I.
(2026). \emph{Care for growth: Making industrial jobs work for women in
Viet Nam} (Policy Summary Note). World Bank Group \& Australian Aid.

Bausch, A., \& Krist, M. (2007). The effect of context-related
moderators on the internationalization-performance relationship:
Evidence from meta-analysis. \emph{Management International Review,
47}(3), 319--347. https://doi.org/10.1007/s11575-007-0019-z

Chiao, Y.-C., Yang, K.-P., \& Yu, C.-M. J. (2006). Performance,
internationalization, and firm-specific advantages of SMEs in a
newly-industrialized economy. \emph{Small Business Economics, 26}(5),
475--492. https://doi.org/10.1007/s11187-005-5599-z

Cameron, A. C., Gelbach, J. B., \& Miller, D. L. (2008). Bootstrap-based
improvements for inference with clustered errors. \emph{The Review of
Economics and Statistics, 90}(3), 414--427.
https://doi.org/10.1162/rest.90.3.414

Cho, H., Kim, C., \& Han, S. (2023). Business group affiliation and the
internationalization--performance relationship: Evidence from Korean
firms. \emph{Journal of International Business Studies, 54}(3),
487--509.

Cohen, W. M., \& Levinthal, D. A. (1990). Absorptive capacity: A new
perspective on learning and innovation. \emph{Administrative Science
Quarterly, 35}(1), 128--152. https://doi.org/10.2307/2393553

Contractor, F. J., Kundu, S. K., \& Hsu, C.-C. (2003). A three-stage
theory of international expansion: The link between multinationality and
performance in the service sector. \emph{Journal of International
Business Studies, 34}(1), 5--18.
https://doi.org/10.1057/palgrave.jibs.8400003

Contractor, F. J., Kumar, V., \& Kundu, S. K. (2007). Nature of the
relationship between international expansion and performance: The case
of emerging market firms. \emph{Journal of World Business, 42}(4),
401--417. https://doi.org/10.1016/j.jwb.2007.06.003

Cuervo-Cazurra, A., Ciravegna, L., Melgarejo, M., \& Lopez, L. (2018).
Home country uncertainty and the internationalization-performance
relationship: Building an uncertainty management capability.
\emph{Journal of World Business, 53}(2), 209--221.
https://doi.org/10.1016/j.jwb.2017.11.003

Dawson, J. F. (2014). Moderation in management research: What, why,
when, and how. \emph{Journal of Business and Psychology, 29}(1), 1--19.
https://doi.org/10.1007/s10869-013-9308-7

Authors. (2024). Details omitted for double-blind review.

Gomes, L., \& Ramaswamy, K. (1999). An empirical examination of the form
of the relationship between multinationality and performance.
\emph{Journal of International Business Studies, 30}(1), 173--187.
https://doi.org/10.1057/palgrave.jibs.8490065

Haans, R. F. J., Pieters, C., \& He, Z.-L. (2016). Thinking about U:
Theorizing and testing U- and inverted U-shaped relationships in
strategy research. \emph{Strategic Management Journal, 37}(7),
1177--1195. https://doi.org/10.1002/smj.2399

Hambrick, D. C. (2007). Upper echelons theory: An update. \emph{Academy
of Management Review, 32}(2), 334--343.
https://doi.org/10.5465/amr.2007.24345254

Hambrick, D. C., \& Mason, P. A. (1984). Upper echelons: The
organization as a reflection of its top managers. \emph{Academy of
Management Review, 9}(2), 193--206.
https://doi.org/10.5465/amr.1984.4277628

International Finance Corporation. (2022, December). \emph{Mind the
gaps: Women in leadership in Viet Nam's banking sector}. IFC, World Bank
Group.

Johanson, J., \& Vahlne, J.-E. (1977). The internationalization process
of the firm---A model of knowledge development and increasing foreign
market commitments. \emph{Journal of International Business Studies,
8}(1), 23--32. https://doi.org/10.1057/palgrave.jibs.8490676

Johanson, J., \& Vahlne, J.-E. (2009). The Uppsala internationalization
process model revisited: From liability of foreignness to liability of
outsidership. \emph{Journal of International Business Studies, 40}(9),
1411--1431. https://doi.org/10.1057/jibs.2009.24

Kirca, A. H., Roth, K., Hult, G. T. M., \& Cavusgil, S. T. (2012). The
role of context in the multinationality--performance relationship: A
meta-analytic review. \emph{Global Strategy Journal, 2}(2), 108--121.
https://doi.org/10.1002/gsj.1028

Lind, J. T., \& Mehlum, H. (2010). With or without U? The appropriate
test for a U-shaped relationship. \emph{Oxford Bulletin of Economics and
Statistics, 72}(1), 109--118.
https://doi.org/10.1111/j.1468-0084.2009.00569.x

Lu, J. W., \& Beamish, P. W. (2004). International diversification and
firm performance: The S-curve hypothesis. \emph{Academy of Management
Journal, 47}(4), 598--609. https://doi.org/10.2307/20159604

Marano, V., Tashman, P., \& Kostova, T. (2016). Escaping the iron cage:
Liabilities of origin and CSR reporting of emerging market multinational
enterprises. \emph{Journal of International Business Studies, 47}(3),
386--408. https://doi.org/10.1057/jibs.2016.7

Mondal, A., Gupta, S., \& Ray, S. (2022). Family ownership, governance,
and internationalization: Evidence from India. \emph{Journal of Business
Research, 138}, 423--434.

North, D. C. (1990). \emph{Institutions, institutional change and
economic performance}. Cambridge University Press.

Pangarkar, N. (2008). Internationalization and performance of small- and
medium-sized enterprises. \emph{Journal of World Business, 43}(4),
475--485. https://doi.org/10.1016/j.jwb.2007.11.009

Richard, O. C., Kirby, S. L., \& Chadwick, K. (2019). The impact of
racial and gender diversity in management on firm performance: An
assessment of different conceptual and measurement approaches.
\emph{Journal of Applied Social Psychology, 39}(7), 1571--1599.

Scott, W. R. (2008). \emph{Institutions and organizations: Ideas and
interests} (3rd ed.). Sage.

Stallkamp, M., \& Schotter, A. P. J. (2021). Platforms without borders?
The international strategies of digital platform firms. \emph{Global
Strategy Journal, 11}(1), 58--80. https://doi.org/10.1002/gsj.1336

UNCTAD. (2023). \emph{World investment report 2023: Investing in
sustainable energy for all}. United Nations.

Verhoef, P. C., Broekhuizen, T., Bart, Y., Bhattacharya, A., Dong, J.
Q., Fabian, N., \& Haenlein, M. (2021). Digital transformation: A
multidisciplinary reflection and research agenda. \emph{Journal of
Business Research, 122}, 889--901.
https://doi.org/10.1016/j.jbusres.2019.09.022

Wernerfelt, B. (1984). A resource-based view of the firm.
\emph{Strategic Management Journal, 5}(2), 171--180.
https://doi.org/10.1002/smj.4250050207

World Bank. (2025). \emph{Enterprise surveys}.
https://www.enterprisesurveys.org

World Bank. (2026, April). \emph{Viet Nam macro monitoring: April 2026}.
World Bank Group, Fiscal Policy and Growth Viet Nam Team.

\begin{center}\rule{0.5\linewidth}{0.5pt}\end{center}

\hypertarget{appendix-a-model-sequence}{%
\subsection{Appendix A: Model
Sequence}\label{appendix-a-model-sequence}}

\textbf{Table 1: Variable definitions and measurement}

\begin{longtable}[]{@{}
  >{\raggedright\arraybackslash}p{(\columnwidth - 6\tabcolsep) * \real{0.2500}}
  >{\raggedright\arraybackslash}p{(\columnwidth - 6\tabcolsep) * \real{0.2500}}
  >{\raggedright\arraybackslash}p{(\columnwidth - 6\tabcolsep) * \real{0.2500}}
  >{\raggedright\arraybackslash}p{(\columnwidth - 6\tabcolsep) * \real{0.2500}}@{}}
\toprule\noalign{}
\begin{minipage}[b]{\linewidth}\raggedright
Variable
\end{minipage} & \begin{minipage}[b]{\linewidth}\raggedright
Definition
\end{minipage} & \begin{minipage}[b]{\linewidth}\raggedright
Source items
\end{minipage} & \begin{minipage}[b]{\linewidth}\raggedright
Coverage
\end{minipage} \\
\midrule\noalign{}
\endhead
\bottomrule\noalign{}
\endlastfoot
lp\_z & Log labour productivity ln(d2/l1), winsorised 1/99, z-scored
within economy-year & d2; l1 & \textasciitilde85\% \\
FSTS & Foreign sales share (direct + indirect exports) & (d3b + d3c) /
100 & \textasciitilde83\% \\
TCI\_z & Tech capability index (quality certification + foreign
technology), z-scored within economy-year & b8, e6 &
\textasciitilde82\% \\
DAI & Digital adoption (website presence, Tier-1 binary) & c22b &
\textasciitilde83\% \\
mgr\_experience & Top manager years in sector & b7 &
\textasciitilde84\% \\
mgr\_female & Top manager is female (binary) & b7a, b6a &
\textasciitilde82\% \\
female\_owner & Majority female ownership (binary) & b4 &
\textasciitilde99\% \\
foreign\_own\_pct & Foreign ownership \% & b6a / 100 &
\textasciitilde52\% \\
firm\_age & Years since establishment & current yr − b5 &
\textasciitilde64\% \\
ln\_size & Log permanent workers & l1 & \textasciitilde84\% \\
ICRV\_group & Institutional regime group (1--6) & ICRV classification &
100\% \\
\end{longtable}

\textbf{Table 2: Model fit and turning points} (50-economy canonical
frame; two-way FE economy+year; CRV1 by economy)

\begin{longtable}[]{@{}
  >{\raggedright\arraybackslash}p{(\columnwidth - 8\tabcolsep) * \real{0.2000}}
  >{\raggedleft\arraybackslash}p{(\columnwidth - 8\tabcolsep) * \real{0.2000}}
  >{\raggedleft\arraybackslash}p{(\columnwidth - 8\tabcolsep) * \real{0.2000}}
  >{\raggedleft\arraybackslash}p{(\columnwidth - 8\tabcolsep) * \real{0.2000}}
  >{\raggedright\arraybackslash}p{(\columnwidth - 8\tabcolsep) * \real{0.2000}}@{}}
\toprule\noalign{}
\begin{minipage}[b]{\linewidth}\raggedright
Model
\end{minipage} & \begin{minipage}[b]{\linewidth}\raggedleft
N
\end{minipage} & \begin{minipage}[b]{\linewidth}\raggedleft
Turning point (\%)
\end{minipage} & \begin{minipage}[b]{\linewidth}\raggedleft
LM p
\end{minipage} & \begin{minipage}[b]{\linewidth}\raggedright
Key additions
\end{minipage} \\
\midrule\noalign{}
\endhead
\bottomrule\noalign{}
\endlastfoot
M1 & 81,022 & --- & --- & FSTS linear (β = +0.219**) \\
M2 & 81,022 & \textbf{51.5} & \textless.001 & FSTS + FSTS² (anchor) \\
M3 & 78,970 & \textbf{43.8} & \textless.001 & M2 + ownership, age, size,
female-owner \\
M4 & 79,080 & \textbf{43.6} & \textless.001 & M3-block + TCI + DAI
(headline) \\
M5 & 79,080 & 44.0 & \textless.001 & M4 + FSTS×TCI, FSTS²×TCI \\
M6 & 79,080 & 47.1 & \textless.001 & M4 + FSTS×DAI, FSTS²×DAI \\
M7 & 75,029 & 44.0 & \textless.001 & M4 + manager experience, female
manager \\
M8 & 79,080 & (73.6) & n.s. & M4 + FSTS×ICRV, FSTS²×ICRV (interactions
NS) \\
\end{longtable}

\emph{Note.} Turning point from −β₁/(2β₂) + mean(FSTS), confirmed by the
Lind--Mehlum (2010) test (LM p). All models carry two-way (economy,
year) fixed effects throughout; SEs clustered by economy. In M8 the ICRV
main effect is absorbed by economy fixed effects and the FSTS×ICRV
interactions are not significant, so the pooled turning point is not
identified; the institutional structure appears instead in the per-ICRV
turning points (Table 4). Reproducible via
\texttt{scripts/p7\_full\_ladder.py}.

\textbf{Table 3: Key coefficient estimates}

\begin{longtable}[]{@{}
  >{\raggedright\arraybackslash}p{(\columnwidth - 12\tabcolsep) * \real{0.1429}}
  >{\raggedleft\arraybackslash}p{(\columnwidth - 12\tabcolsep) * \real{0.1429}}
  >{\raggedleft\arraybackslash}p{(\columnwidth - 12\tabcolsep) * \real{0.1429}}
  >{\raggedleft\arraybackslash}p{(\columnwidth - 12\tabcolsep) * \real{0.1429}}
  >{\raggedleft\arraybackslash}p{(\columnwidth - 12\tabcolsep) * \real{0.1429}}
  >{\raggedleft\arraybackslash}p{(\columnwidth - 12\tabcolsep) * \real{0.1429}}
  >{\raggedleft\arraybackslash}p{(\columnwidth - 12\tabcolsep) * \real{0.1429}}@{}}
\toprule\noalign{}
\begin{minipage}[b]{\linewidth}\raggedright
Term
\end{minipage} & \begin{minipage}[b]{\linewidth}\raggedleft
M2
\end{minipage} & \begin{minipage}[b]{\linewidth}\raggedleft
M4
\end{minipage} & \begin{minipage}[b]{\linewidth}\raggedleft
M5
\end{minipage} & \begin{minipage}[b]{\linewidth}\raggedleft
M6
\end{minipage} & \begin{minipage}[b]{\linewidth}\raggedleft
M7
\end{minipage} & \begin{minipage}[b]{\linewidth}\raggedleft
M8
\end{minipage} \\
\midrule\noalign{}
\endhead
\bottomrule\noalign{}
\endlastfoot
FSTS (β₁) & +1.188*** & +0.499*** & +0.457*** & +0.704*** & +0.505*** &
+0.806† \\
FSTS² (β₂) & −1.398*** & −0.721*** & −0.652** & −0.924*** & −0.720*** &
−0.624 \\
TCI\_z & --- & +0.108*** & +0.115*** & +0.108*** & +0.108*** &
+0.107*** \\
DAI & --- & +0.219*** & +0.219*** & +0.201*** & +0.218*** & +0.220*** \\
FSTS×TCI & --- & --- & +0.068 & --- & --- & --- \\
FSTS²×TCI & --- & --- & −0.119 & --- & --- & --- \\
FSTS×DAI & --- & --- & --- & −0.270 & --- & --- \\
FSTS²×DAI & --- & --- & --- & +0.252 & --- & --- \\
mgr experience (z) & --- & --- & --- & --- & +0.002 & --- \\
female manager & --- & --- & --- & --- & −0.104*** & --- \\
FSTS×ICRV & --- & --- & --- & --- & --- & −0.086 \\
FSTS²×ICRV & --- & --- & --- & --- & --- & −0.013 \\
N & 81,022 & 79,080 & 79,080 & 79,080 & 75,029 & 79,080 \\
\end{longtable}

\emph{Note.} Two-way FE (economy, year); SEs clustered by economy. *** p
\textless{} .001, ** p \textless{} .01, * p \textless{} .05, † p
\textless{} .10. Capability (TCI, DAI) carries robust positive level
effects throughout; none of the FSTS×capability or FSTS×ICRV interaction
terms is significant; the pooled curve is reshaped by neither capability
nor a linear institutional gradient. Reproducible via
\texttt{scripts/p7\_full\_ladder.py}.

\textbf{Table 4: Turning points within each ICRV regime}
(capability-adjusted M4 form; two-way FE)

\begin{longtable}[]{@{}
  >{\raggedright\arraybackslash}p{(\columnwidth - 10\tabcolsep) * \real{0.1667}}
  >{\raggedleft\arraybackslash}p{(\columnwidth - 10\tabcolsep) * \real{0.1667}}
  >{\raggedleft\arraybackslash}p{(\columnwidth - 10\tabcolsep) * \real{0.1667}}
  >{\raggedleft\arraybackslash}p{(\columnwidth - 10\tabcolsep) * \real{0.1667}}
  >{\raggedleft\arraybackslash}p{(\columnwidth - 10\tabcolsep) * \real{0.1667}}
  >{\raggedright\arraybackslash}p{(\columnwidth - 10\tabcolsep) * \real{0.1667}}@{}}
\toprule\noalign{}
\begin{minipage}[b]{\linewidth}\raggedright
ICRV group
\end{minipage} & \begin{minipage}[b]{\linewidth}\raggedleft
N
\end{minipage} & \begin{minipage}[b]{\linewidth}\raggedleft
β₁
\end{minipage} & \begin{minipage}[b]{\linewidth}\raggedleft
β₂
\end{minipage} & \begin{minipage}[b]{\linewidth}\raggedleft
Turning point (\%)
\end{minipage} & \begin{minipage}[b]{\linewidth}\raggedright
Reading
\end{minipage} \\
\midrule\noalign{}
\endhead
\bottomrule\noalign{}
\endlastfoot
I Advanced innovation (incl.~Japan) & 5,581 & +0.707 & −0.504 & ≈ 80 &
near-linear (beyond range) \\
II Advanced resource (GCC) & 2,075 & +0.774 & −0.730 & 56.6 & broad
inverted-U \\
III Upper-middle & 12,055 & +0.174 & −0.189 & 56.4 & weak/flat \\
IV Lower-middle transition & 42,094 & +0.689 & −1.012 & \textbf{≈ 43} &
sharp inverted-U \\
V Emerging & 15,457 & +0.235 & −0.455 & 36.8 & dissolving \\
VI Pacific SIDS & 1,818 & −0.681 & +0.870 & --- & no inverted-U (FIP) \\
\end{longtable}

\emph{Note.} Within-group quadratic with controls and two-way FE. The
three-zone structure (sharp in the transition regime, near-linear at the
top, dissolving at the bottom) is the paper's central institutional
finding. Consistent with the dissertation's canonical per-ICRV values.

\begin{center}\rule{0.5\linewidth}{0.5pt}\end{center}

\hypertarget{appendix-b-tfp-production-function-estimates-robustness-reference}{%
\subsection{Appendix B: TFP Production Function Estimates (Robustness
Reference)}\label{appendix-b-tfp-production-function-estimates-robustness-reference}}

Sector-level production function coefficients used to construct
alternative TFP-based dependent variables for robustness checks. Two
specifications are available: VA=f(K,L) (VAKL, value-added formulation)
and Y=f(K,L,M) (YKLM, gross-output formulation), each estimated
separately for 20 two-digit ISIC manufacturing sectors with high-income
economy interaction terms and year fixed effects (2009--2025).
Coefficients and t-statistics are stored in
\texttt{replication/Annex\_TFP\_regression\_tables\_March\_11\_2026.xlsx}.
TFP residuals from these regressions can be used as an alternative
dependent variable to ln(LP) to test whether the inverted-U relationship
is robust to output-per-worker vs.~total-factor productivity
operationalisations.

\begin{center}\rule{0.5\linewidth}{0.5pt}\end{center}

\end{document}
